% appendix/styleguide/styleguide.tex
% mainfile: ../../perfbook.tex
% SPDX-License-Identifier: CC-BY-SA-3.0

\chapter{风格指南}
\label{chp:app:styleguide:Style Guide}
%
\Epigraph{De gustibus non est disputandum.（趣味无争辩。）}{拉丁格言}

本附录是风格指南的集合，旨在作为改善perfbook一致性的参考。
它还包含若干建议及其实验示例。

\Cref{sec:app:styleguide:Paul's Conventions} 描述了基本的标点和拼写规则。
\Cref{sec:app:styleguide:NIST Style Guide} 解释了与单位符号相关的规则。
\Cref{sec:app:styleguide:LaTeX Conventions} 总结了\LaTeX 特有的约定。

\section{Paul的约定}
\label{sec:app:styleguide:Paul's Conventions}

以下是Paul约定的列表，由他对Akira关于perfbook标点策略的提问的回答汇编而成。

\begin{itemize}
\item （关于标点符号和引号）
  尽管我自己是美国人，但对于这类书籍，英式方法更好，
  因为它消除了如下歧义：
  \begin{quote}
    Type ``\nbco{ls -a},'' look for the file ``\co{.},''
    and file a bug if you don't see it.
  \end{quote}

  以下写法则清晰得多：
  \begin{quote}
    Type ``\nbco{ls -a}'', look for the file ``\co{.}'',
    and file a bug if you don't see it.
  \end{quote}
\item 美式英语拼写：用``color''而非``colour''。
\item 牛津逗号：用``a, b, and~c''而非``a, b and~c''。
  这是一种任意约定。
  当牛津逗号导致歧义时，应改写句子，
  例如引入编号：``a, b, and c and~d''应写成
  ``(1)~a, (2)~b, and (3)~c and~d''。
\item 用斜体表示强调。
  应谨慎使用。
\item 标识符使用 \verb|\co{}|，URL使用 \verb|\url{}|，
  文件名使用 \verb|\path{}|。
\item 日期应使用无歧义的格式。
  不要用``mm/dd/yy''或``dd/mm/yy''，而应使用``July 26, 2016''
  或``26 July 2016''或``26-Jul-2016''或``2016/07/26''。
  例如，我倾向于在文件名中使用 \path{yyyy.mm.ddA}。
\item 北美关于句号和缩写的规则。
  例如，以下两个句子都不会被合理地解读为两个独立的句子：
  \begin{itemize}
  \item Say hello, to Mr.~Jones.
  \item If it looks like she sprained her ankle, call Dr.~Smith and
    then tell her to keep the ankle iced and elevated.
  \end{itemize}

  一个有歧义的例子：
  \begin{quote}
    If I take the cow, the pig, the horse, etc. George will be upset.
  \end{quote}
  可以用更多词语来改写：
  \begin{quote}
    If I take the cow, the pig, the horse, or much of anything else,
    George will be upset.
  \end{quote}
  或者：
  \begin{quote}
    If I take the cow, the pig, the horse, etc., George will be upset.
  \end{quote}
\item 我不喜欢在标题中使用 \& 符号（``\&''），但如果使用它
  能避免标题中出现换行，我有时也会使用。
\item 提及词语时，我使用引号。
  引入新词时，我使用 \verb|\emph{}|。
\end{itemize}

以下是关于\LaTeX\ 源代码中标点符号的约定。

\begin{itemize}
\item 在冒号（\co{:}）和句末之后换行。
  这样可以避免单空格/双空格的争论，同时还有一个好处，
  即能更清晰地显示长段落中单个句子的修改。
\end{itemize}

\section{NIST风格指南}
\label{sec:app:styleguide:NIST Style Guide}

\subsection{单位符号}
\label{sec:app:styleguide:Unit Symbol}

\subsubsection{SI单位符号}
\label{sec:app:styleguide:SI Unit Symbol}

NIST风格指南~\cite[Chapter 5]{NIST:SP:330:2019}
规定了以下规则（为perfbook重新措辞）。

\begin{itemize}
\item 当SI单位符号如``ns''、``MHz''和``K''（开尔文）
用在数值之后时，数值和符号之间应放置窄空格。

在 \LaTeX{} 中，窄空格可以用 \qco{\\,} 序列来编写。
例如，
\begin{quote}
  ``2.4\,GHz'', rather then ``2.4GHz''.
\end{quote}

\item 即使数值用作形容词修饰时，也应放置窄空格。
  例如，
\begin{quote}
  ``a~10\,ms interval'', rather than ``a~10\=/ms interval'' nor
  ``a~10ms interval''.
\end{quote}
\end{itemize}

微（\micro :$10^{-6}$）的符号可以借助``gensymb'' \LaTeX\ 宏包
轻松排版。
宏 \qco{\\micro} 可以在文本模式和数学模式中使用。
要排版“微秒”的符号，可以使用
\qco{\\micro s}。
例如，
\begin{quote}
  10\,\micro s
\end{quote}

注意，数学模式的 \qco{\\mu} 默认是斜体的，不应
用作前缀。
一个不正确的例子：
\begin{quote}
  10\,$\mu $s (math mode \qco{\\mu})
\end{quote}

\subsubsection{非SI单位符号}
\label{sec:app:styleguide:Non-SI Unit Symbol}

虽然NIST风格指南不涵盖非SI单位符号如"KB"、"MB"和"GB"，但应遵循相同的规则。

例如：

\begin{quote}
  ``A~240\,GB hard drive'', rather than ``a~240\=/GB hard drive''
  nor ``a~240GB hard drive''.
\end{quote}

严格来说，NIST指南要求我们使用二进制前缀
``Ki''、``Mi''或``Gi''来表示~$2^{10}$的幂。
但是，我们接受JEDEC约定，使用``K''、``M''
和``G''作为描述内存容量时的二进制前缀~\cite{JEDEC:dict:prefixmega}。

一个可接受的例子：
\begin{quote}
  ``8\,GB of main memory'', meaning ``8\,GiB of main memory''.
\end{quote}

此外，也可以仅使用``K''、``M''或``G''作为缩写
附在数值之后，例如``4K~entries''。
在这种情况下，缩写前不需要空格。
例如，

\begin{quote}
  ``8K entries'', rather than ``8\,K entries''.
\end{quote}

如果在中间放置空格，该符号看起来就像单位符号，
容易造成混淆。
注意，``K''和``k''分别代表$2^{10}$和$10^3$。
``M''可以代表$2^{20}$或$10^6$，``G''可以代表
$2^{30}$或$10^9$。
在讨论大致数量级时，这些歧义不会造成混淆。

\subsubsection{度数符号}
\label{sec:app:styleguide:Degree Symbol}

角度符号（\degree）前面不需要任何空格。
NIST风格指南明确这样规定。

度数符号也可以借助gensymb
宏包轻松排版。
宏 \qco{\\degree} 可以在文本模式和数学模式中使用。

例如：

\begin{quote}
  $45\degree$, rather than $45\,\degree$.
\end{quote}

\subsubsection{百分号}
\label{sec:app:styleguide:Percent Symbol}

NIST风格指南将百分号（\%）视为与SI单位
符号相同。

\begin{quote}
  50\pct\ possibility, rather than 50\% possibility.
\end{quote}

\subsubsection{字体样式}
\label{sec:app:styleguide:Font Style}

引用自NIST检查列表~\cite[\#6]{NIST:UnitCheckList}：

\begin{quote}
  Variables and quantity symbols are in italic type.
  Unit symbols are in roman type.
  Numbers should generally be written in roman type.
  These rules apply irrespective of the typeface used in the surrounding
  text.
\end{quote}

例如，
\begin{quote}
  {\textit e}（基本电荷）
\end{quote}

另一方面，数学常数如自然对数的底数
应使用正体~\cite{NIST:TypeFaces}。
例如，

\begin{quote}
  $\mathrm{e}^x$
\end{quote}

%\footnote{
%  See \url{https://tex.stackexchange.com/questions/119248/}
%  for the historical reason.}

\subsection{尚未遵循的NIST指南}
\label{sec:app:styleguide:NIST Guides Yet To Be Followed}

有少数情况下未遵循NIST风格指南。
在这些情况下，遵循的是其他英语惯例。

\subsubsection{数字分组}
\label{sec:app:styleguide:Digit Grouping}

引用自NIST检查列表~\cite[\#16]{NIST:UnitCheckList}：

\begin{quote}
  The digits of numerical values having more than four digits on either
  side of the decimal marker are separated into groups of three using
  a thin, fixed space counting from both the left and right of the decimal
  marker.
  Commas are not used to separate digits into groups of three.
\end{quote}

\begin{quote}
\begin{tabular}{ll}
  NIST Example:& 15\,739.012\,53\,ms\\
  Our convention:& 15,739.01253\,ms\\
\end{tabular}
\end{quote}

在 \LaTeX\ 编码中，按照NIST指南的建议放置窄空格是很繁琐的。
``siunitx''宏包提供的 \verb|\num{}| 命令可以
帮助我们遵循这一规则。
它还可以帮助我们克服不同的约定差异。
我们可以在序言中选择特定的数字分组样式作为默认值，
或者为每个 \verb|\num{}| 命令指定选项，
如在\cref{tab:app:styleguide:Digit-Grouping Style}中所示。

\newcommand{\NumDigitGrpA}{12 345}
\newcommand{\NumDigitGrpB}{12.345}
\newcommand{\NumDigitGrp}{1 234 567.89}
\begin{table}
\small\centering
\begin{tabular}{lrrr}\toprule
  Style & \multicolumn{3}{c}{Outputs of \texttt{\textbackslash num\{\}}} \\
  \midrule
  NIST/SI (English) & \num[group-separator={\,},group-digits=integer]{\NumDigitGrpA} &
    \num[group-separator={\,},group-digits=integer]{\NumDigitGrpB} &
      \num[group-separator={\,},group-digits=integer]{\NumDigitGrp} \\
  SI (French) & \num[locale=FR,group-separator={\,}]{\NumDigitGrpA} &
    \num[locale=FR,group-separator={\,}]{\NumDigitGrpB} &
      \num[locale=FR,group-separator={\,}]{\NumDigitGrp} \\
  English & \num[group-separator={,},group-digits=integer]{\NumDigitGrpA} &
    \num[group-separator={,},group-digits=integer]{\NumDigitGrpB} &
      \num[group-separator={,},group-digits=integer]{\NumDigitGrp} \\
  French & \num[locale=FR,group-separator={\,}]{\NumDigitGrpA} &
    \num[locale=FR,group-separator={\,}]{\NumDigitGrpB} &
      \num[locale=FR,group-separator={\,}]{\NumDigitGrp} \\
  Other Europe & \num[group-separator={.},output-decimal-marker={,},group-digits=integer]{\NumDigitGrpA} &
    \num[group-separator={.},output-decimal-marker={,},group-digits=integer]{\NumDigitGrpB} &
      \num[group-separator={.},output-decimal-marker={,},group-digits=integer]{\NumDigitGrp} \\
\bottomrule
\end{tabular}
\caption{数字分组样式}
\label{tab:app:styleguide:Digit-Grouping Style}
\end{table}

如\cref{tab:app:styleguide:Digit-Grouping Style}中所示，
将句号和逗号用作小数点以外的用途会造成混淆，
应该避免，尤其是在面向全球读者的文档中。

通过使用 \verb|\num{}| 命令来标记常量小数值，
\LaTeX\ 源代码可以不受任何特定约定的约束。

由于其开源策略，这种方法应该能为perfbook提供
更好的“可移植性”。

\section{\LaTeX\ 约定}
\label{sec:app:styleguide:LaTeX Conventions}

美观的\LaTeX\ 文档需要进一步考虑字体样式的正确使用、换行例外等。
本节总结了\LaTeX\ 特有的准则。

\subsection{等宽字体}
\label{sec:app:styleguide:Monospace Font}

等宽字体（或打字机字体）在本教材中被大量使用。
perfbook中关于等宽字体的首要策略是避免直接使用 \qco{\\texttt} 或 \qco{\\tt} 宏。
强烈建议使用能够表明你需要该字体原因的宏或环境。

本节解释了这些宏和环境的用例。

\subsubsection{代码片段}
\label{sec:app:styleguide:Code Snippet}

由于 \qco{verbatim} 环境是一种较为原始的代码列表包含方式，
我们已经转向使用 \qco{fancyvrb} 包来处理代码片段的方案。

该方案的目标是直接从 \path{CodeSamples} 目录下的代码示例中
提取代码片段的\LaTeX\ 源码。
它还使得在代码示例中嵌入行标签成为可能，
这些标签可以在\LaTeX\ 源码中被引用。
这减轻了保持正文中行号与代码片段中行号一致的负担。

代码片段的提取由几个 perl 脚本和 Makefile 中的构建规则来处理。
我们利用 \co{fancyvrb} 包的转义特性
将行标签嵌入为注释。

我们过去使用 \qco{verbatimbox} 包提供的 \qco{verbbox} 环境。
\Cref{sec:app:styleguide:Code Snippet (Obsolete)} 描述了
\co{verbbox} 如何自动生成行号，
但这些行号无法在\LaTeX\ 源码中被引用。

让我们先来看看当前方案中代码片段是如何编写的。
\qco{Verbatim} 有三个自定义环境。
\qco{VerbatimL} 用于 \qco{listing} 环境内的浮动代码片段。
\qco{VerbatimN} 用于启用行计数的内联代码片段。
\qco{VerbatimU} 用于不带行计数的内联代码片段。
它们在导言区中的定义如下所示：

\begin{VerbatimU}
\DefineVerbatimEnvironment{VerbatimL}{Verbatim}%
  {fontsize=\scriptsize,numbers=left,numbersep=5pt,%
    xleftmargin=9pt,obeytabs=true,tabsize=2}
\AfterEndEnvironment{VerbatimL}{\vspace*{-9pt}}
\DefineVerbatimEnvironment{VerbatimN}{Verbatim}%
  {fontsize=\scriptsize,numbers=left,numbersep=3pt,%
    xleftmargin=5pt,xrightmargin=5pt,obeytabs=true,%
    tabsize=2,frame=single}
\DefineVerbatimEnvironment{VerbatimU}{Verbatim}%
  {fontsize=\scriptsize,numbers=none,xleftmargin=5pt,%
    xrightmargin=5pt,obeytabs=true,tabsize=2,%
    samepage=true,frame=single}
\end{VerbatimU}

% Another option would be the ``lstlisting'' environment provided
%  by the ``listings'' package. We are already using its ``lstinline''
%  command in the definition of \co{\\co\{\}} macro.

\begin{listing}
\fvset{fontsize=\scriptsize,numbers=left,numbersep=5pt,xleftmargin=9pt,obeytabs=true,tabsize=8,commandchars=\%\~\^}
\begin{fcvlabel}[ln:app:styleguide:LaTeX Source of Sample Code Snippet (Current)]
\VerbatimInput{appendix/styleguide/samplecodesnippetfcv.tex}
\end{fcvlabel}
\vspace*{-9pt}
\caption{示例代码片段的 \LaTeX\ 源码（当前方案）}
\label{lst:app:styleguide:LaTeX Source of Sample Code Snippet (Current)}
\end{listing}

\begin{listing}
\begin{fcvlabel}[ln:base1]			%lnlbl~beg:fcvlabel^
\begin{VerbatimL}[commandchars=\$\[\]]
/*
 * Sample Code Snippet
 */
#include <stdio.h>
int main(void)
{
	printf("Hello world!\n");	$lnlbl[printf]
	return 0;			$lnlbl[return]
}
\end{VerbatimL}
\end{fcvlabel}					%lnlbl~end:fcvlabel^
\caption{示例代码片段}
\label{lst:app:styleguide:Sample Code Snippet}
\end{listing}


一个代码片段示例的\LaTeX\ 源码显示在
\cref{lst:app:styleguide:LaTeX Source of Sample Code Snippet (Current)}
中，其排版效果显示在
\cref{lst:app:styleguide:Sample Code Snippet} 中。

行标签通过 \qco{$lnlbl[]} 命令来指定。
\co{VarbatimL} 环境的 \qco{commandchars} 选项所指定的字符
被 \co{fancyvrb} 包用于将 \qco{$lnlbl[]} 替换为
\qco{\\lnlbl\{\}}。
这些字符的选取应确保它们不会出现在代码片段的其他位置。

\cref{lst:app:styleguide:Sample Code Snippet} 中的标签
\qco{printf} 和 \qco{return}
可以按如下方式引用：

\begin{VerbatimU}
\begin{fcvref}[ln:base1]
\Clnref{printf, return} can be referred
to from text.
\end{fcvref}
\end{VerbatimU}

上述代码的排版结果如下段所示：

\begin{quote}
\begin{fcvref}[ln:base1]
\Clnref{printf, return} can be referred
to from text.
\end{fcvref}
\end{quote}

宏 \qco{\\lnlbl\{\}} 和 \qco{\\lnref\{\}} 在
导言区中的定义如下：

\begin{VerbatimU}
\newcommand{\lnlblbase}{}
\newcommand{\lnlbl}[1]{%
  \phantomsection\label{\lnlblbase:#1}}
\newcommand{\lnrefbase}{}
\newcommand{\lnref}[1]{\ref{\lnrefbase:#1}}
\end{VerbatimU}

环境 \qco{fcvlabel} 和 \qco{fcvref} 的定义如下所示：

\begin{VerbatimU}
\newenvironment{fcvlabel}[1][]{%
  \renewcommand{\lnlblbase}{#1}%
  \ignorespaces}{\ignorespacesafterend}
\newenvironment{fcvref}[1][]{%
  \renewcommand{\lnrefbase}{#1}%
  \ignorespaces}{\ignorespacesafterend}
\end{VerbatimU}

\begin{fcvref}[ln:app:styleguide:LaTeX Source of Sample Code Snippet (Current)]
\cref{lst:app:styleguide:LaTeX Source of Sample Code Snippet (Current)}
中\clnrefrange{beg:fcvlabel}{end:fcvlabel} 所示的\LaTeX\ 源码的主体部分
可以通过 perl 脚本 \path{utilities/fcvextract.pl}
从\cref{lst:app:styleguide:Source of Code Sample} 的代码示例中提取。
所有相关的提取规则都以构建规则的形式记录在
顶层 \path{Makefile} 和一个依赖生成脚本
(\path{utilities/gen_snippet_d.pl}) 中。
\end{fcvref}

\begin{listing*}
\fvset{fontsize=\scriptsize,numbers=left,numbersep=5pt,xleftmargin=9pt,obeytabs=true,tabsize=8}
\VerbatimInput{appendix/styleguide/samplecodesnippet.c}
\vspace*{-9pt}
\caption{带 ``snippet'' 元命令的代码示例源码}
\label{lst:app:styleguide:Source of Code Sample}
\end{listing*}

如你所见，\cref{lst:app:styleguide:Source of Code Sample}
在 C（C++ 风格）注释中包含了元命令。
这些元命令由 \path{utilities/fcvextract.pl} 解释，
该脚本通过代码示例的文件扩展名来区分注释风格的类型。

可在代码示例中使用的元命令如下所列：

\begin{itemize}[noitemsep]
\item \texttt{\textbackslash begin\{snippet\}[<options>]}
\item \texttt{\textbackslash end\{snippet\}}
\item \texttt{\textbackslash lnlbl\{<label string>\}}
\item \texttt{\textbackslash fcvexclude}
\item \texttt{\textbackslash fcvblank}
\end{itemize}

\texttt{\textbackslash begin\{snippet\}} 元命令的 \qco{<options>}
是一个以逗号分隔的选项列表，如下所示：

\begin{itemize}[noitemsep]
\item \co{labelbase=<label base string>}
\item \co{keepcomment=yes}
\item \co{gobbleblank=yes}
\item \co{commandchars=\\X\\Y\\Z}
\end{itemize}

\qco{labelbase} 选项是必需的，传给它的字符串
将被传递给
``\texttt{\textbackslash begin\{fcvlabel\}[<label base string>]}'' 命令，如
\cref{lst:app:styleguide:LaTeX Source of Sample Code Snippet (Current)} 的
\clnrefr{ln:app:styleguide:LaTeX Source of Sample Code Snippet (Current):beg:fcvlabel} 所示。
\qco{keepcomment=yes} 选项告诉 \co{fcvextract.pl} 保留注释块。
否则，C 源代码中的注释块将被省略。
\qco{gobbleblank=yes} 选项将移除生成的代码片段中的空行或空白行。
\qco{commandchars} 选项将原样传递给 \co{VerbatimL} 环境。
目前，该选项也是必需的，且必须位于上述选项列表的末尾。
其他类型的选项（如果有的话）也会被传递给 \co{VerbatimL}
环境。

\qco{\\lnlbl} 命令在处理过程中会被转换以反映
转义字符的选择。\footnote{
	\qco{\\lnlbl} 命令周围构成注释的字符
	也会被删除，与 \qco{keepcomment} 设置无关。
}
包含 \qco{\\fcvexclude} 的源代码行将被移除。
当指定了 \qco{gobbleblank=yes} 选项时，
可以使用 \qco{\\fcvblank} 来保留空行。

可以有多对 \texttt{\textbackslash begin\{snippet\}} 和 \texttt{\textbackslash end\{snippet\}}，
只要它们具有唯一的 \qco{labelbase} 字符串即可。

我们的 \qco{labelbase} 唯一命名空间的命名方案如下：

\begin{VerbatimU}
ln:<Chapter/Subdirectory>:<File Name>:<Function Name>
\end{VerbatimU}

由 \qco{herdtools7} 命令（如 \qco{litmus7} 和 \qco{herd7}）
处理的 litmus 测试在此方案中存在问题。
这些命令对注释可以放置的位置以及注释中允许的字符有特殊的规则。
它们还禁止某些标记出现在注释中。
（注释中的标记听起来可能很奇怪，但它们确实有这样的限制。）

例如，litmus 测试中的第一个标记必须是
\qco{C}、\qco{PPC}、\qco{X86}、\qco{LISA} 等之一，
用于指示测试的类型。
这意味着 litmus 测试的开头不允许有注释。

类似地，\qco{exists}、\qco{filter}
和~\qco{locations} 等几个标记表示 litmus 测试主体的结束。
一旦它们出现在 litmus 测试中，注释应该采用
OCaml 风格（\qco{(* ... *)}）。
这些标记即使出现在注释中也保持相同的含义！

字符对 \qco{\{} 和 \qco{\}} 在 C 类型的测试中
也有特殊含义。
它们用于分隔 litmus 测试中的各个部分。

第一对 \qco{\{} 和 \qco{\}} 包围的是初始化部分。
这部分中的注释也应该采用 OCaml 形式。

在 litmus 测试的注释中也不能使用 \qco{\{} 和 \qco{\}}。

litmus 测试中不允许使用的注释示例如下所示：

\begin{fcvlabel}[ln:app:styleguide:Bad comments in Litmus Test]
\begin{VerbatimN}[tabsize=8]
// Comment at first
C C-sample
// Comment with { and } characters
{
x=2;  // C style comment in initialization
}

P0(int *x}
{
	int r1;

	r1 = READ_ONCE(*x);  // Comment with "exists"
}

[...]

exists (0:r1=0)  // C++ style comment after test body
\end{VerbatimN}
\end{fcvlabel}

为避免解析错误，litmus 测试（C 类型）中的元命令按以下方式嵌入。

\begin{fcvlabel}[ln:app:styleguide:Sample Source of Litmus Test]
\begin{VerbatimN}[tabsize=8]
C C-SB+o-o+o-o
//\begin[snippet][labelbase=ln:base,commandchars=\%\@\$]

{
1:r2=0			(*\lnlbl[initr2]*)
}

P0(int *x0, int *x1)		//\lnlbl[P0:b]
{
	int r2;

	WRITE_ONCE(*x0, 2);
	r2 = READ_ONCE(*x1);
}				//\lnlbl[P0:e]

P1(int *x0, int *x1)
{
	int r2;

	WRITE_ONCE(*x1, 2);
	r2 = READ_ONCE(*x0);
}

//\end[snippet]
exists (1:r2=0 /\ 0:r2=0)  (* \lnlbl[exists_] *)
\end{VerbatimN}
\end{fcvlabel}

上面的示例通过脚本 \path{utilities/reorder_ltms.pl}
被转换为以下中间代码。\footnote{
	目前仅支持 C 类型的 litmus 测试。
}
该中间代码可以由通用脚本 \path{utilities/fcvextract.pl} 处理。

\begin{fcvlabel}[ln:app:styleguide:Intermediate Source of Litmus Test]
\begin{VerbatimN}[tabsize=8]
// Do not edit!
// Generated by utillities/reorder_ltms.pl
//\begin{snippet}[labelbase=ln:base,commandchars=\%\@\$]
C C-SB+o-o+o-o

{
1:r2=0			//\lnlbl{initr2}
}

P0(int *x0, int *x1)		//\lnlbl{P0:b}
{
	int r2;

	WRITE_ONCE(*x0, 2);
	r2 = READ_ONCE(*x1);
}				//\lnlbl{P0:e}

P1(int *x0, int *x1)
{
	int r2;

	WRITE_ONCE(*x1, 2);
	r2 = READ_ONCE(*x0);
}

exists (1:r2=0 /\ 0:r2=0)  \lnlbl{exists_}
//\end{snippet}
\end{VerbatimN}
\end{fcvlabel}

请注意，由于注释的限制，每个 litmus 测试的源文件最多只能包含一对
\texttt{\textbackslash begin[snippet]} 和 \texttt{\textbackslash end[snippet]}。

\subsubsection{代码片段（已废弃）}
\label{sec:app:styleguide:Code Snippet (Obsolete)}

使用 \qco{verbatimbox} 包编写的代码片段的
\LaTeX\ 源码示例显示在
\cref{lst:app:styleguide:LaTeX Source of Sample Code Snippet (Obsolete)}
中，其排版效果显示在
\cref{lst:app:styleguide:Sample Code Snippet (Obsolete)} 中。

\begin{listing}
\begin{fcvlabel}[ln:app:styleguide:samplecodesnippetlstlbl]
\fvset{fontsize=\scriptsize,numbers=left,numbersep=5pt,xleftmargin=9pt,commandchars=\%\@\$}
\VerbatimInput{appendix/styleguide/samplecodesnippetlstlbl.tex}
\end{fcvlabel}
\vspace*{-9pt}
\caption{示例代码片段的 \LaTeX\ 源码（已废弃）}
\label{lst:app:styleguide:LaTeX Source of Sample Code Snippet (Obsolete)}
\end{listing}

\begin{listing}
{ \scriptsize
\begin{verbbox}[\LstLineNo]
/*
 * Sample Code Snippet
 */
#include <stdio.h>
int main(void)
{
  printf("Hello world!\n");
  return 0;
}
\end{verbbox}
}
\centering
\theverbbox
\caption{示例代码片段（已废弃）}
\label{lst:app:styleguide:Sample Code Snippet (Obsolete)}
\end{listing}


\co{verbbox} 的自动编号功能通过
在 verbbox 选项中指定的 ``\verb|\LstLineNo|'' 宏来启用
（\cref{lst:app:styleguide:LaTeX Source of Sample Code Snippet (Obsolete)} 的
\clnrefr{ln:app:styleguide:samplecodesnippetlstlbl:lineno}）。
该宏在 \path{perfbook.tex} 的导言区中定义如下：

\begin{VerbatimU}
\newcommand{\LstLineNo}
  {\makebox[5ex][r]{\arabic{VerbboxLineNo}\hspace{2ex}}}
\end{VerbatimU}

\qco{verbatim} 环境用于行数过多而无法放入一栏的代码列表。
它还用于避免过多的浮动对象使\LaTeX\ 不堪重负。
这些代码列表正在被转换为使用 \co{VerbatimN} 环境的方案。

\subsubsection{标识符}
\label{sec:app:styleguide:Identifier}

我们使用 ``\verb|\co{}|'' 宏来内联标识符。
（``co'' 代表 ``code''。）

将标识符放入 \verb|\co{}| 后，其名称中的下划线字符
无需在 \LaTeX\ 源代码中进行转义。
这也便于在源文件中搜索它们。
此外，\verb|\co{}| 宏还具有在特定字符序列处允许换行的功能。
当前定义允许在下划线（\tco{_}）、
两个连续的下划线（\tco{__}）、空格或
运算符 \tco{->} 处换行。

\subsubsection{表格和标题中的标识符}
\label{sec:app:styleguide:Identifier inside Table and Heading}

虽然 \verb|\co{}| 命令在正文中内联使用很方便，
但由于其换行功能，它是脆弱的。
当它在 \qco{tabular} 环境或其派生环境（如 \qco{tabularx}）中使用时，
会干扰这些环境的列宽估算。
此外，\verb|\co{}| 不能安全地用于章节标题或
描述标题中。

作为变通方案，我们在表格和标题中使用 ``\verb|\tco{}|'' 命令。
它不具备在特定序列处换行的功能，但
仍然免去了转义下划线的麻烦。

在正文中使用时，\verb|\tco{}| 允许在空格处换行。

\subsubsection{等宽字体的其他用例}
\label{sec:app:styleguide:Other Use Cases of Monospace Font}

对于 URL，我们使用 \qco{hyperref} 包提供的 ``\verb|\url{}|'' 命令。
它会生成指向 URL 的超链接。

对于路径名，我们使用 ``\verb|\path{}|'' 命令。
它不会生成超链接。

\verb|\url{}| 和 \verb|\path{}| 都允许在
\qco{/}、\qco{-} 和 \qco{.} 处换行。\footnote{
  如果 URL 或路径名包含长串不可断字符，溢出可能成为一个问题。
}

对于不希望被换行的短等宽字体语句，我们使用
``\verb|\nbco{}|''（non-breakable co）宏。

\subsubsection{限制}
\label{sec:app:styleguide:Limitations}

本节介绍的宏在某些情况下可能无法按预期工作。
\Cref{tab:app:styleguide:Limitation of Monospace Macro}
列出了这些限制。

\begin{table}
\renewcommand*{\arraystretch}{1.2}\centering\footnotesize
\begin{tabular}{@{}lll@{}}\toprule
  Macro &  Need Escape & Should Avoid \\
  \midrule
  \texttt{\textbackslash co}, \texttt{\textbackslash nbco} & \texttt{\textbackslash}, \%, \{, \} & \\
  \texttt{\textbackslash tco}  & \# & \%, \{, \}, \texttt{\textbackslash} \\
  \bottomrule
\end{tabular}
\caption{等宽字体宏的限制}
\label{tab:app:styleguide:Limitation of Monospace Macro}
\end{table}

虽然 \verb|\co{}| 要求某些字符必须转义，
但它可以包含任何字符。

另一方面，\verb|\tco{}| 无法正确处理
\qco{\%}、\qco{\{}、\qco{\}} 和 \qco{\\}。
如果使用 \qco{\\} 对它们进行转义，
转义字符会出现在最终结果中。
如果你需要对包含大量需要转义字符的字符串使用等宽字体，
可以在正文中使用 \qco{\\verb} 命令。\footnote{
  不过 \texttt{\textbackslash verb} 命令也不是万能的。
  例如，你不能在脚注中使用它。
  如果这样做，你会看到一个致命的 \LaTeX\ 错误。
  一个变通方案是使用 \co{fancyvrb} 包提供的
  名为 \texttt{\textbackslash VerbatimFootnotes} 的宏。
  不幸的是，perfbook 由于与 \co{footnotebackref} 包的冲突
  而无法使用它。
  }

\subsection{交叉引用}
\label{sec:app:styleguide:Cross-Reference}

对\namecrefs{chp:Introduction}、
\namecrefs{sec:intro:Parallel Programming Goals}、
\namecrefs{lst:app:styleguide:Source of Code Sample}等的交叉引用，
以往通过名称与裸 \verb|\ref{}|
命令的组合来表示，方式如下：

\begin{VerbatimN}
Chapter~\ref{chp:Introduction},
Table~\ref{tab:app:styleguide:Digit-Grouping Style}
\end{VerbatimN}

这是传统的交叉引用方式。
然而，在每个交叉引用上手动添加名称既繁琐又容易出错。
\co{cleveref} 包提供了一种更好的交叉引用方式。
以下是一些示例：

\begin{VerbatimN}
\Cref{chp:Introduction},
\cref{sec:intro:Parallel Programming Goals},
\cref{chp:app:styleguide:Style Guide},
\cref{tab:app:styleguide:Digit-Grouping Style}, and
\cref{lst:app:styleguide:Source of Code Sample} are
examples of cross\-/references.
\end{VerbatimN}

上述代码的排版效果如下：

\begin{quote}
\Cref{chp:Introduction},
\cref{sec:intro:Parallel Programming Goals},
\cref{chp:app:styleguide:Style Guide},
\cref{tab:app:styleguide:Digit-Grouping Style}, and
\cref{lst:app:styleguide:Source of Code Sample} are
examples of cross\-/references.
\end{quote}

如你所见，交叉引用的命名是自动化的。
当前设置对 \verb|\Cref{}| 和 \verb|\cref{}| 都生成首字母大写的名称，
但前者应在句首使用。

我们正在进行向 \co{cleveref} 风格交叉引用的转换。

对代码片段行号的交叉引用可以通过
\verb|\Clnref{}| 和 \verb|\clnref{}| 宏以类似方式完成，
它们模仿了 \co{cleveref} 的行为。
前者使用 ``Line'' 作为引用名称，
后者使用 ``line''。

请参阅 \co{cleveref} 的文档以了解更多关于其智能功能的信息。

\subsection{不可断空格}
\label{sec:app:styleguide:Non Breakable Spaces}

在 \LaTeX\ 的惯例中，强烈建议正确使用不可断空格。
它们可以防止单个数字或短变量名成为孤儿或寡妇，
否则会导致文本乍看之下令人困惑。

前面提到的放在单位符号前面的窄空格是不可断的。

其他使用不可断空格（\LaTeX\ 源代码中的 ``\verb|~|''，
通常称为 ``nbsp''）的场景包括以下几种（不完全列举）。

\begin{itemize}
\item 引用章节：
  \begin{quote}
    Please refer to \cref{sec:app:styleguide:NIST Style Guide}.
  \end{quote}
\item 指明 CPU 编号或 Thread 名称：
  \begin{quote}
    After they load the pointer, CPUs~1 and~2 will see the stored
    value.
  \end{quote}
\item 短变量名：
  \begin{quote}
    The results will be stored in variables~\co{a} and~\co{b}.
  \end{quote}
\end{itemize}

\subsection{连字符与破折号}
\label{sec:app:styleguide:Hyphenation and Dashes}

\subsubsection{复合词中的连字符}
\label{sec:app:styleguide:Hyphenation in Compound Word}

在普通 \LaTeX 中，像 ``high-frequency'' 这样的复合词
只能在连字符处断行。
这有时会导致排版效果不佳。
例如：

\begin{center}\begin{minipage}{2.6in}\vspace{0.6\baselineskip}
  High-frequency radio wave, high-frequency radio wave,
  high-frequency radio wave, high-frequency radio wave,
  high-frequency radio wave, high-frequency radio wave.
\vspace{0.6\baselineskip}\end{minipage}\end{center}

通过使用 ``extdash'' 包提供的快捷方式 \qco{\\-/}，
perfbook 中启用了复合词各组成部分的断词功能。\footnote{
  作为启用该快捷方式的代价，我们无法使用普通
  \LaTeX 的快捷方式 \qco{\\-} 来指定断词点。
  请使用 \path{pfhyphex.tex} 来添加此类例外。
}

使用 \qco{\\-/} 的示例：

\begin{center}\begin{minipage}{2.6in}\vspace{0.6\baselineskip}
  High\-/frequency radio wave, high\-/frequency radio wave,
  high\-/frequency radio wave, high\-/frequency radio wave,
  high\-/frequency radio wave, high\-/frequency radio wave.
\vspace{0.6\baselineskip}\end{minipage}\end{center}

\subsubsection{不可断连字符}
\label{sec:app:styleguide:Non Breakable Hyphen}

我们希望像 ``x\=/coordinate''、
``y\=/coordinate'' 等带连字符的复合术语不要在单个字母后的连字符处断行。

要使连字符不可断，可以使用同样由 ``extdash'' 包提供的
快捷方式 \qco{\\=/}。

不使用快捷方式的示例：

\begin{center}\begin{minipage}{2.55in}\vspace{0.6\baselineskip}
x-, y-, and z-coordinates; x-, y-, and z-coordinates;
x-, y-, and z-coordinates; x-, y-, and z-coordinates;
x-, y-, and z-coordinates; x-, y-, and z-coordinates;
\vspace{0.6\baselineskip}\end{minipage}\end{center}

使用 \qco{\\-/} 的示例：

\begin{center}\begin{minipage}{2.55in}\vspace{0.6\baselineskip}
x-, y-, and z\-/coordinates; x-, y-, and z\-/coordinates;
x-, y-, and z\-/coordinates; x-, y-, and z\-/coordinates;
x-, y-, and z\-/coordinates; x-, y-, and z\-/coordinates;
\vspace{0.6\baselineskip}\end{minipage}\end{center}

使用 \qco{\\=/} 的示例：

\begin{center}\begin{minipage}{2.55in}\vspace{0.6\baselineskip}
x-, y-, and z\=/coordinates; x-, y-, and z\=/coordinates;
x-, y-, and z\=/coordinates; x-, y-, and z\=/coordinates;
x-, y-, and z\=/coordinates; x-, y-, and z\=/coordinates;
\vspace{0.6\baselineskip}\end{minipage}\end{center}

注意，\qco{\\=/} 与 \qco{\\-/} 一样，
也会启用复合词各组成部分的断词功能。

\subsubsection{长破折号}
\label{sec:app:styleguide:Em Dash}

长破折号用于表示插入语。
在 perfbook 中，长破折号两边不加空格。
在 \LaTeX\ 源代码中，长破折号用 \qco{---} 表示。

示例（引自\cref{sec:app:whymb:Cache Structure}）：
\begin{quote}
  This disparity in speed---more than two orders of magnitude---has
  resulted in the multi-megabyte caches found on modern CPUs.
\end{quote}

\subsubsection{短破折号}
\label{sec:app;styleguide:En Dash}

在 \LaTeX\ 惯例中，短破折号（\==）用于表示（主要是）数字的范围。
perfbook 过去的版本没有遵循这一规则，而是使用
普通连字符（\=/）来表示此类情况。

现在 \texttt{\textbackslash clnrefrange}、\texttt{\textbackslash crefrange}
及其变体（它们会生成短破折号）已被用于交叉引用的范围，
其余几十个其他类型范围的简单连字符也已统一转换为短破折号，
以保持一致性。

使用简单连字符的示例：

\begin{quote}
\begin{fcvref}[ln:app:styleguide:samplecodesnippetlstlbl]
  Lines~\lnref{b}\=/\lnref{e} in
  \cref{lst:app:styleguide:LaTeX Source of Sample Code Snippet (Obsolete)}
  are the contents of the verbbox environment.
  The box is output by the \texttt{\textbackslash theverbbox} macro on \clnref{theverbbox}.
\end{fcvref}
\end{quote}

使用短破折号的示例：

\begin{quote}
\begin{fcvref}[ln:app:styleguide:samplecodesnippetlstlbl]
  Lines~\lnref{b}\==\lnref{e} in
  \cref{lst:app:styleguide:LaTeX Source of Sample Code Snippet (Obsolete)}
  are the contents of the verbbox environment.
  The box is output by the \texttt{\textbackslash theverbbox} macro on \clnref{theverbbox}.
\end{fcvref}
\end{quote}

\subsubsection{数字减号}
\label{sec:app:styleguide:Numerical Minus Sign}

数字减号应编码为数学模式的减号，
即 \qco{$-$}。\footnote{此规则假设数学模式使用与文本模式相同的正体字形。
  我们默认的字体选择满足这一假设。
\IfSansSerif{
  截至 2021 年 8 月，使用无衬线字体的实验性目标（如
  ``1csf'' 和 ``ebsf''）在数学模式中\emph{确实}使用了不同的字体。}{}
}
例如，

\begin{quote}
  $-30$, rather than -30.
\end{quote}

\subsection{标点符号}
\label{sec:app:styleguide:Punctuation}

\subsubsection{省略号}
\label{sec:app:styleguide:Ellipsis}

在等宽字体中，省略号可以用一系列句点来表示。
例如：

\begin{quote}
  \verb|Great ... So how do I fix it?|
\end{quote}

然而，在比例字体中，一系列句点会以紧凑的间距打印，如下所示：

\begin{quote}
  Great ... So how do I fix it?
\end{quote}

标准 \LaTeX\ 为此目的定义了 \verb|\dots| 宏。
但它在间距均匀性方面存在一个缺陷。
``ellipsis'' 包重新定义了 \verb|\dots| 宏以修复
此问题。\footnote{确切地说，被重新定义的是 \texttt{\textbackslash textellipsis} 宏。
  \texttt{\textbackslash dots} 宏在数学模式下的行为不受影响。
  ``amsmath'' 包有另一个 \texttt{\textbackslash dots} 的定义，
  目前在 perfbook 中未使用。}
使用 \verb|\dots| 后，上述示例的排版效果如下：

\begin{quote}
  Great \dots So how do I fix it?
\end{quote}

注意，为 ``ellipsis'' 包指定的 ``xspace'' 选项
会根据省略号后面的内容来调整其后的间距。

例如：

\begin{itemize}[itemsep=.2ex]
\item He said, ``I~\dots really don't remember~\dots''
\item Sequence A: (one, two, three, \dots)
\item Sequence B: (4, 5, \dots, $n$)
\end{itemize}

如你所见，逗号前面放置了额外的间距。

\verb|\dots| 宏也可以在数学模式中使用：

\begin{itemize}[itemsep=.2ex]
\item Sequence C: $(1, 2, 3, 5, 8, \dots)$
\item Sequence D: $(10, 12, \dots, 20)$
\end{itemize}

\verb|\ldots| 宏的行为与 \verb|\dots| 宏相同。

\subsubsection{句号}
\label{sec:app:styleguide:Full Stop}

\LaTeX\ 默认将空格前的句号视为句末，
并放置稍宽的间距（双倍间距）。
此规则有一个例外，即当句号紧跟在大写字母后面时，
\LaTeX\ 会认为它表示缩写，
从而放置正常间距。

为了让 \LaTeX\ 使用正确的间距，需要标注此类例外。
例如，给定以下 \LaTeX\ 源代码：

\begin{VerbatimU}
\begin{quote}
	Lock~1 is owned by CPU~A.
	Lock~2 is owned by CPU~B.  (Bad.)

	Lock~1 is owned by CPU~A\@.
	Lock~2 is owned by CPU~B\@.  (Good.)
\end{quote}
\end{VerbatimU}
%
输出结果如下：

\begin{quote}
	Lock~1 is owned by CPU~A.
	Lock~2 is owned by CPU~B.  (Bad.)

	Lock~1 is owned by CPU~A\@.
	Lock~2 is owned by CPU~B\@.  (Good.)
\end{quote}

另一方面，当句号跟在小写字母后面时，
例如 ``Mr.~Smith'' 中的情况，如果没有正确提示，
输出中会出现较宽的间距。
这种提示可以通过以下几种方式来完成。

给定以下源代码，

\begin{VerbatimU}
\begin{itemize}[nosep]
	\item Mr. Smith (bad)
	\item Mr.~Smith (good)
	\item Mr.\ Smith (good)
	\item Mr.\@ Smith (good)
\end{itemize}
\end{VerbatimU}

\noindent%
结果如下所示：

\begin{itemize}[nosep]
	\item Mr. Smith (bad)
	\item Mr.~Smith (good)
	\item Mr.\ Smith (good)
	\item Mr.\@ Smith (good)
\end{itemize}

\subsection{浮动对象格式}
\label{sec:app:styleguide:Floating Object Format}

\subsubsection{表格中的线条}
\label{sec:app:styleguide:Ruled Line in Table}

据说使用纯 \LaTeX\ 线条绘制的表格看起来很丑。\footnote{
  \url{https://www.inf.ethz.ch/personal/markusp/teaching/guides/guide-tables.pdf}
}
应避免使用竖线，横线应谨慎使用，尤其是在结构简单的表格中。

\IfTblCpTop{}{
\floatstyle{plaintop}
\restylefloat{table}
\setlength{\abovetopsep}{-2pt}
\addtolength{\abovecaptionskip}{-2.5pt}
}

\newcommand{\TLo}{T_\mathrm{L}}
\newcommand{\THi}{T_\mathrm{H}}
\newcommand{\CPf}{C_\mathrm{P}}

\Cref{tab:app:styleguide:Refrigeration Power Consumption}
（对应于一个已删除章节中的表格）
是使用 ``booktabs'' 和 ``xcolor'' 包的功能绘制的。
请注意，booktabs 的横线不能与表格中的竖线混合使用。\footnote{
  还有一个名为 ``arydshln'' 的包，可以提供在表格中使用的虚线。
  一些实验性的示例在
  \cref{sec:app:styleguide:Table Layout Experiment}中展示。
}

\begin{table}
\rowcolors{1}{}{lightgray}
\renewcommand*{\arraystretch}{1.2}\centering\small
\begin{tabular}{lrrr}\toprule
Situation
	& $T$ (K)
		& $\CPf$ & \parbox[b]{.75in}{\raggedleft Power per watt\par waste heat (W)} \\
\midrule
Dry Ice
	& $195$
		& $1.990$
			& 0.5 \\
Liquid N$_2$
	& $77$
		& $0.356$
			& 2.8 \\
Liquid H$_2$
	& $20$
		& $0.073$
			& 13.7 \\
Liquid He
	& $4$
		& $0.0138$
			& 72.3 \\
IBM~Q	& $0.015$
		& $0.000051$
			& 19,500.0 \\
\bottomrule
\end{tabular}
\caption{制冷功耗}
\label{tab:app:styleguide:Refrigeration Power Consumption}
\end{table}

\subsubsection{标题位置}
\label{sec:app:styleguide:Position of Caption}

在 \LaTeX\ 的惯例中，表格的标题通常放在表格上方。
原因在于你查看表格时的视线流动方向。
大多数表格的顶部有一行标题。
你自然会首先看向表格的顶部。
放在表格底部的标题会打断这种流程。
代码片段也是如此，因为代码是从上往下阅读的。

对于代码片段，列表环境选用的 ``ruled'' 样式会将标题放在顶部。
示例参见\cref{lst:app:styleguide:Sample Code Snippet}。

对于表格，标题的位置通过前导区中的
\verb|\floatstyle{}| 和 \verb|\restylefloat{}| 宏来调整。

通过给 \verb|\abovecaptionskip| 变量设置较小的值，
可以减少标题下方的垂直间距，
但这也会影响图表的标题。

在使用 ``booktabs'' 包横线的表格中，
通过给 \texttt{\textbackslash abovetopsep} 变量设置负值，
可以进一步减少标题与表格之间的垂直间距，
该变量控制 \texttt{\textbackslash toprule} 的行为。

\subsection{改进候选}
\label{sec:app:styleguide:Improvement Candidates}

\begin{figure*}\centering
\ebresizewidth{
\begin{minipage}[t][][t]{2.1in}
\resizebox{2.1in}{!}{\includegraphics{cartoons/1kHz}}
\caption{1\,kHz 定时器轮}
\label{fig:app:styleguide:Timer Wheel at 1kHz}
\end{minipage}
\qquad
\begin{minipage}[t][][t]{2.3in}
\resizebox{2.3in}{!}{\includegraphics{cartoons/100kHz}}
\caption{100\,kHz 定时器轮}
\label{fig:app:styleguide:Timer Wheel at 100kHz}
\end{minipage}
}
\end{figure*}

\begin{listing*}%
\caption{消息传递石蕊测试（使用 subfig）}%
\label{lst:app:styleguide:Message-Passing Litmus Test (subfig)}%
{\scriptsize%
\begin{verbbox}[\LstLineNo]
C C-MP+o-wmb-o+o-o.litmus

{
}

P0(int* x0, int* x1) {

  WRITE_ONCE(*x0, 2);
  smp_wmb();
  WRITE_ONCE(*x1, 2);

}

P1(int* x0, int* x1) {

  int r2;
  int r3;

  r2 = READ_ONCE(*x1);
  r3 = READ_ONCE(*x0);

}


exists (1:r2=2 /\ 1:r3=0)
\end{verbbox}
}
\centering
\hspace*{\fill}
\subfloat[未强制排序]{
  \theverbbox
  \label{sublst:app:styleguide:Not Enforcing Order}
}
\hspace{\fill}
{\scriptsize%
\begin{verbbox}[\LstLineNo]
C C-MP+o-wmb-o+o-rmb-o.litmus

{
}

P0(int* x0, int* x1) {

  WRITE_ONCE(*x0, 2);
  smp_wmb();
  WRITE_ONCE(*x1, 2);

}

P1(int* x0, int* x1) {

  int r2;
  int r3;

  r2 = READ_ONCE(*x1);
  smp_rmb();
  r3 = READ_ONCE(*x0);

}

exists (1:r2=2 /\ 1:r3=0)
\end{verbbox}
}%
\subfloat[强制排序]{%
  \theverbbox
  \label{sublst:app:styleguide:Enforcing Order}
}\hspace*{\fill}%
\end{listing*}

在perfbook中还有一些领域有待尝试，这些改进将进一步改善其外观。
本节列出了这些候选项。

\subsubsection{分组相关图表/代码清单}
\label{sec:app:styleguide:Grouping Related Figures/Listings}

为了防止一对紧密相关的图表或代码清单被放置在不同的页面上，
最好将它们组合到一个浮动对象中。
``subfig'' 包提供了实现此功能的特性。\footnote{
  组合图表的一个问题可能是 \LaTeX\ 源代码的复杂性。}

两个浮动对象可以通过 \texttt{\textbackslash parbox} 或
\co{minipage} 并排放置。
例如，
\cref{fig:advsync:Timer Wheel at 1kHz,fig:advsync:Timer Wheel at 100kHz}
可以使用一对 \co{minipage} 组合在一起，
如\cref{fig:app:styleguide:Timer Wheel at 1kHz,%
fig:app:styleguide:Timer Wheel at 100kHz}所示。

使用 subfig 包，
\cref{lst:memorder:Message-Passing Litmus Test (No Ordering),%
lst:memorder:Enforcing Order of Message-Passing Litmus Test}
可以组合在一起，如
\cref{lst:app:styleguide:Message-Passing Litmus Test (subfig)}
所示，并带有子标题（空行有微小改动）。

请注意，它们不能像
\cref{fig:app:styleguide:Timer Wheel at 1kHz,%
fig:app:styleguide:Timer Wheel at 100kHz}
那样组合，因为 ``ruled'' 样式会阻止其标题被正确排版。

子标题可以通过组合 \verb|\cref{}| 宏和
\verb|\subref{}| 宏来引用，例如
``\cref{lst:app:styleguide:Message-Passing Litmus Test (subfig)}\,%
\subref{sublst:app:styleguide:Not Enforcing Order}''。

也可以通过 \verb|\cref{}| 宏来引用，例如
``\cref{sublst:app:styleguide:Enforcing Order}''。
请注意生成格式的差异。
为使 \verb|\cref{}| 引用正常工作，需要将组合浮动对象的
\verb|\label{}| 宏放在子浮动体定义之前。
否则，生成的标题编号会与实际编号相差一。

\subsubsection{表格布局实验}
\label{sec:app:styleguide:Table Layout Experiment}

本节展示了一些使用booktabs、xcolors和arydshln包的实验性表格。
正文中对应的表格已使用此处展示的某种格式进行了转换。
本节的源代码可以作为在正文中添加新表格时的参考。

\begin{table}
\rowcolors{1}{}{lightgray}
\renewcommand*{\arraystretch}{1.1}
\sisetup{group-minimum-digits=4,group-separator={,}}
\centering\small
\begin{tabular}
  {
    l
    S[table-format = 9.1]
    S[table-format = 9.1]
  }
	\toprule
	Operation		& \multicolumn{1}{r}{Cost (ns)}
			& {\parbox[b]{.7in}{\raggedleft Ratio\\(cost/clock)}} \\
	\midrule
	Clock period		&           0.5	&           1.0 \\
	Best-case CAS		&           7.0	&          14.6 \\
	Best-case lock		&          15.4	&          32.3 \\
	Blind CAS		&           7.2	&          15.2 \\
	CAS			&          18.0	&          37.7 \\
	Blind CAS (off-core)	&          47.5	&          99.8 \\
	CAS (off-core)		&         101.9	&         214.0 \\
	Blind CAS (off-socket)	&         148.8	&         312.5 \\
	CAS (off-socket)	&         442.9	&         930.1 \\
	Comms Fabric		&       5 000	&      10 500	\\
	Global Comms		& 195 000 000	& 409 500 000   \\
	\bottomrule
\end{tabular}
\caption{CPU 0 视角的同步机制（8 路 Intel Xeon Platinum 8176 CPU @ 2.10GHz 系统）}
\label{tab:app:styleguide:CPU 0 View of Synchronization Mechanisms on 8-Socket System With Intel Xeon Platinum 8176 CPUs @ 2.10GHz}
\end{table}

在\cref{tab:app:styleguide:CPU 0 View of Synchronization Mechanisms on 8-Socket System With Intel Xeon Platinum 8176 CPUs @ 2.10GHz}
（对应于\cref{tab:cpu:CPU 0 View of Synchronization Mechanisms on 8-Socket System With Intel Xeon Platinum 8176 CPUs at 2.10GHz}）中，
使用了 ``siunitx'' 包提供的 ``S'' 列说明符来对齐数字。

\Cref{tab:app:styleguide:Synchronization and Reference Counting}
（对应于\cref{tab:together:Synchronizing Reference Counting}）
是具有复杂表头的表格示例。
在\cref{tab:app:styleguide:Synchronization and Reference Counting}中，
中线的间隔对应于转换前用双竖线表示的区分。
附加的框线图例解释了矩阵中使用的缩写。
此处两种类型的内存屏障用下标表示。
\cref{tab:together:Synchronizing Reference Counting}中
没有这些图例和下标，因为在那里它们是多余的。

\begin{table}
\renewcommand*{\arraystretch}{1.25}
\rowcolors{3}{}{lightgray}
\small
\centering
\begin{tabular}{lcccc}
	\toprule
	& \multicolumn{4}{c}{Release} \\
	\cmidrule(l){2-5}
	Acquisition	& Locks
				& \parbox[c]{.5in}{Reference\\Counts}
					& \parbox[c]{.5in}{Hazard\\Pointers}
						& RCU \\
	\cmidrule{1-1} \cmidrule(l){2-5}
	Locks		& $-$	& CAM\textsubscript{R}	& M	& CA  \\
	\parbox[c][6ex]{.6in}{Reference\\Counts}
			& A	& AM\textsubscript{R}    & M	& A   \\
	\parbox[c][6ex]{.6in}{Hazard\\Pointers}
			& M	& M	& M	& M   \\
	RCU		& CA	& M\textsubscript{A}CA	& M	& CA  \\
	\bottomrule
\end{tabular}

\vspace{5pt}\hfill
\framebox[\width]{\footnotesize\setlength{\tabcolsep}{3pt}
\rowcolors{1}{}{}
  \begin{tabular}{lrp{2in}}
	Key:	& A: & Atomic counting \\
		& C: & Check combined with the atomic acquisition operation \\
		& M: & Full memory barriers required \\
		& M\textsubscript{R}: & Memory barriers required only on release \\
		& M\textsubscript{A}: & Memory barriers required on acquire \\
  \end{tabular}
}
\caption{同步与引用计数}
\label{tab:app:styleguide:Synchronization and Reference Counting}
\end{table}

\Cref{tab:app:styleguide:Cache Coherence Example}
（对应于\cref{tab:app:whymb:Cache Coherence Example}）
是以表格形式绘制的序列图。

\begin{table*}
\small
\centering
\renewcommand*{\arraystretch}{1.2}
\rowcolors{6}{}{lightgray}
% "6" is chosen due to disturbance of row count by cmidrule.
% The command definition is:
%     \rowcolors{<row>}{<odd-row color>}{<even-row color>}
% Here, <row> specifies the row count where the coloring start.
% In this table, the "Seq = 0" row is the 3rd row, so a "3" would
% be a right choice.
% However, because of the \cmidrule{} commands used in the heading,
% internal row count of the "Seq = 0" row becomes "6".
% This is why the 3rd row has the background color of <even-row color>.
%
% \cline of plain LaTeX also interfares the row count.
\ebresizewidth{
\begin{tabular}{rclcccccc}
	\toprule
	& & & \multicolumn{4}{c}{CPU Cache} & \multicolumn{2}{c}{Memory} \\
	\cmidrule(lr){4-7} \cmidrule(l){8-9}
	Sequence \# & CPU \# & Operation & 0 & 1 & 2 & 3 & 0 & 8 \\
	\cmidrule(r){1-3} \cmidrule(lr){4-7} \cmidrule(l){8-9}
%	Seq CPU Operation	------------- CPU -------------   - Memory -
%				   0	   1	   2	   3	    0   8
	0   &   & Initial State	& $-$/I & $-$/I & $-$/I & $-$/I   & V & V \\
	1   & 0 & Load 0 (cache value miss)
				& 0/S   & $-$/I & $-$/I & $-$/I   & V & V \\
	2   & 2 & Load 0 (cache value miss)
				& 0/S   & $-$/I & 0/S   & $-$/I   & V & V \\
	3.1 & 0 & Load 8 (collision, invalidation)
				& $-$/I & $-$/I & 0/S   & $-$/I   & V & V \\
	3.2 & 0 & Load 8 (cache value miss)
				& 8/S   & $-$/I & 0/S   & $-$/I   & V & V \\
	4   & 3 & Load 0 (cache value miss)
				& 8/S   & $-$/I & 0/S   & 0/S     & V & V \\
	5.1 & 3 & CAS 0 (cache value hit)
				& 8/S   & $-$/I & 0/S   & 0/S     & V & V \\
	5.2 & 3 & CAS 0 (cache permission miss)
				& 8/S   & $-$/I & $-$/I & 0/E     & V & V \\
	5.3 & 3 & CAS 0		& 8/S   & $-$/I & $-$/I & 0/M     & I & V \\
	6.1 & 1 & Store 8 (cache value miss)
				& $-$/I & 8/E   & $-$/I & 0/M     & I & V \\
	6.2 & 1 & Store 8 	& $-$/I & 8/M   & $-$/I & 0/M     & I & I \\
	7   & 3 & CAS 0		& $-$/I & 8/M   & $-$/I & 0/M     & I & I \\
	8.1 & 1 & Load 0 (collision, flush)
				& $-$/I & $-$/I & $-$/I & 0/M     & I & V \\
	8.2 & 1 & Load 0 (cache value miss)
				& $-$/I & 0/S   & $-$/I & 0/S     & I & V \\
	9   & 3 & Load 0 (cache value hit)
				& $-$/I & 0/S   & $-$/I & 0/S     & I & V \\
	\bottomrule
\end{tabular}
}
\caption{缓存一致性示例}
\label{tab:app:styleguide:Cache Coherence Example}
\end{table*}

\Cref{tab:app:styleguide:RCU Publish-Subscribe and Version Maintenance APIs}
是\cref{tab:defer:RCU Publish-Subscribe and Version Maintenance APIs}的调整版本。
此处移除了原表中的 ``Category'' 列，
改为在中线下方以粗体字行来表示类别。
这一改动使得 ``xcolor'' 包的 \verb|\rowcolors{}| 命令更容易正常工作。

\Cref{tab:app:styleguide:RCU Publish-Subscribe and Version Maintenance APIs (colortbl)}
是另一个版本，它保留了原始列，且仅在某类别有多行时才对行着色。
这是通过组合 ``xcolor'' 的 \verb|\rowcolors{}| 和
``colortbl'' 包的 \verb|\cellcolor{}| 命令实现的
（\verb|\cellcolor{}| 会覆盖 \verb|\rowcolors{}|）。

在\cref{tab:defer:RCU Publish-Subscribe and Version Maintenance APIs}中，
为简洁起见选择了后一种不带部分行着色的布局。

\begin{table*}
\rowcolors{2}{}{blue!15}
\renewcommand*{\arraystretch}{1.1}
\footnotesize
\centering
\ebresizewidth{
\begin{tabular}{lll}
\toprule
	Primitives &
		Availability &
			Overhead \\
\midrule
	\multicolumn{3}{l}{\bfseries List traversal} \\
	\tco{list_for_each_entry_rcu()} &
		2.5.59 &
			Simple instructions (memory barrier on Alpha) \\
\midrule
	\multicolumn{3}{l}{\bfseries List update} \\
	\tco{list_add_rcu()} &
		2.5.44 &
			Memory barrier \\
	\rowcolor{lightgray}\tco{list_add_tail_rcu()} &
		2.5.44 &
			Memory barrier \\
	\tco{list_del_rcu()} &
		2.5.44 &
			Simple instructions \\
	\rowcolor{lightgray}\tco{list_replace_rcu()} &
		2.6.9 &
			Memory barrier \\
	\tco{list_splice_init_rcu()} &
		2.6.21 &
			Grace-period latency \\
\midrule
	\multicolumn{3}{l}{\bfseries Hlist traversal} \\
	\tco{hlist_for_each_entry_rcu()} &
		2.6.8 &
			Simple instructions (memory barrier on Alpha) \\
\midrule
	\multicolumn{3}{l}{\bfseries Hlist update} \\
	\tco{hlist_add_after_rcu()} &
		2.6.14 &
			Memory barrier \\
	\rowcolor{lightgray}\tco{hlist_add_before_rcu()} &
		2.6.14 &
			Memory barrier \\
	\tco{hlist_add_head_rcu()} &
		2.5.64 &
			Memory barrier \\
	\rowcolor{lightgray}\tco{hlist_del_rcu()} &
		2.5.64 &
			Simple instructions \\
	\tco{hlist_replace_rcu()} &
		2.6.15 &
			Memory barrier \\
\midrule
	\multicolumn{3}{l}{\bfseries Pointer traversal} \\
	\tco{rcu_dereference()} &
		2.6.9 &
			Simple instructions (memory barrier on Alpha) \\
\midrule
	\multicolumn{3}{l}{\bfseries Pointer update} \\
	\tco{rcu_assign_pointer()} &
		2.6.10 &
			Memory barrier \\
\bottomrule
\end{tabular}
}
\caption{RCU 发布-订阅和版本维护 API}
\label{tab:app:styleguide:RCU Publish-Subscribe and Version Maintenance APIs}
\end{table*}

\begin{table*}
\renewcommand*{\arraystretch}{1.2}
\rowcolors{3}{lightgray}{}
\footnotesize
\centering
\ebresizewidth{
\begin{tabular}{lllp{1.2in}}\toprule
Category &
	Primitives &
		Availability &
			Overhead \\
\midrule
List traversal &
	\tco{list_for_each_entry_rcu()} &
		2.5.59 &
			Simple instructions (memory barrier on Alpha) \\
\midrule
\cellcolor{white}List update &
	\tco{list_add_rcu()} &
		2.5.44 &
			Memory barrier \\
&
	\tco{list_add_tail_rcu()} &
		2.5.44 &
			Memory barrier \\
\cellcolor{white} &
	\tco{list_del_rcu()} &
		2.5.44 &
			Simple instructions \\
&
	\tco{list_replace_rcu()} &
		2.6.9 &
			Memory barrier \\
\cellcolor{white} &
	\tco{list_splice_init_rcu()} &
		2.6.21 &
			Grace-period latency \\
\midrule
Hlist traversal &
	\tco{hlist_for_each_entry_rcu()} &
		2.6.8 &
			Simple instructions (memory barrier on Alpha) \\
\midrule
\cellcolor{white}Hlist update &
	\tco{hlist_add_after_rcu()} &
		2.6.14 &
			Memory barrier \\
&
	\tco{hlist_add_before_rcu()} &
		2.6.14 &
			Memory barrier \\
\cellcolor{white} &
	\tco{hlist_add_head_rcu()} &
		2.5.64 &
			Memory barrier \\
&
	\tco{hlist_del_rcu()} &
		2.5.64 &
			Simple instructions \\
\cellcolor{white} &
	\tco{hlist_replace_rcu()} &
		2.6.15 &
			Memory barrier \\
\midrule\hiderowcolors
Pointer traversal &
	\tco{rcu_dereference()} &
		2.6.9 &
			Simple instructions (memory barrier on Alpha) \\
\midrule
Pointer update &
	\tco{rcu_assign_pointer()} &
		2.6.10 &
			Memory barrier \\
\bottomrule
\end{tabular}
}
\caption{RCU 发布-订阅和版本维护 API}
\label{tab:app:styleguide:RCU Publish-Subscribe and Version Maintenance APIs (colortbl)}
\end{table*}

\Cref{tab:app:styleguide:Memory Misordering: Store-Buffering Sequence of Events}
（对应于\cref{tab:memorder:Memory Misordering: Store-Buffering Sequence of Events}）
也是以表格对象形式绘制的序列图。

\begin{table*}
\rowcolors{6}{}{lightgray}
\renewcommand*{\arraystretch}{1.1}
\small
\centering\OneColumnHSpace{-0.1in}
\ebresizewidth{
\begin{tabular}{rllllll}
	\toprule
	& \multicolumn{3}{c}{CPU 0} & \multicolumn{3}{c}{CPU 1} \\
	\cmidrule(l){2-4} \cmidrule(l){5-7}
	& Instruction & Store Buffer & Cache &
		Instruction & Store Buffer & Cache \\
	\cmidrule{1-1} \cmidrule(l){2-4} \cmidrule(l){5-7}
	1 & (Initial state) & & \tco{x1==0} &
		(Initial state) & & \tco{x0==0} \\
	2 & \tco{x0 = 2;} & \tco{x0==2} & \tco{x1==0} &
		\tco{x1 = 2;} & \tco{x1==2} & \tco{x0==0} \\
	3 & \tco{r2 = x1;} (0) & \tco{x0==2} & \tco{x1==0} &
		\tco{r2 = x0;} (0) & \tco{x1==2} & \tco{x0==0} \\
	4 & (Read-invalidate) & \tco{x0==2} & \tco{x0==0} &
		(Read-invalidate) & \tco{x1==2} & \tco{x1==0} \\
	5 & (Finish store) & & \tco{x0==2} &
		(Finish store) & & \tco{x1==2} \\
	\bottomrule
\end{tabular}
}
\caption{内存乱序：Store-Buffering 事件序列}
\label{tab:app:styleguide:Memory Misordering: Store-Buffering Sequence of Events}
\end{table*}

\Cref{tab:app:styleguide:Refrigeration Power Consumption (arydshln)}
展示了\cref{tab:app:styleguide:Refrigeration Power Consumption}的另一版本，
使用了 arydshln 包的虚线水平线和竖线。

\setlength\dashlinedash{.5pt}
\setlength\dashlinegap{1pt}

\begin{table}
\renewcommand*{\arraystretch}{1.2}\centering\small
\begin{tabular}{l:r:r:r}\toprule
Situation
	& $T$ (K)
		& $\CPf$ & \parbox[b]{.75in}{\raggedleft Power per watt\par waste heat (W)} \\
\hline
Dry Ice
	& $195$
		& $1.990$
			& 0.5 \\ \hdashline
Liquid N$_2$
	& $77$
		& $0.356$
			& 2.8 \\ \hdashline
Liquid H$_2$
	& $20$
		& $0.073$
			& 13.7 \\ \hdashline
Liquid He
	& $4$
		& $0.0138$
			& 72.3 \\ \hdashline
IBM~Q	& $0.015$
		& $0.000051$
			& 19,500.0 \\
\bottomrule
\end{tabular}
\caption{制冷功耗}
\label{tab:app:styleguide:Refrigeration Power Consumption (arydshln)}
\end{table}

在此情况下，竖向虚线似乎是不必要的。
不带竖线的版本展示在
\cref{tab:app:styleguide:Refrigeration Power Consumption (arydshln-2)}中。

\begin{table}
\renewcommand*{\arraystretch}{1.2}\centering\small
\begin{tabular}{lrrr}\toprule
Situation
	& $T$ (K)
		& $\CPf$ & \parbox[b]{.75in}{\raggedleft Power per watt\par waste heat (W)} \\
\midrule
Dry Ice
	& $195$
		& $1.990$
			& 0.5 \\ \hdashline
Liquid N$_2$
	& $77$
		& $0.356$
			& 2.8 \\ \hdashline
Liquid H$_2$
	& $20$
		& $0.073$
			& 13.7 \\ \hdashline
Liquid He
	& $4$
		& $0.0138$
			& 72.3 \\ \hdashline
IBM~Q	& $0.015$
		& $0.000051$
			& 19,500.0 \\
\bottomrule
\end{tabular}
\caption{制冷功耗}
\label{tab:app:styleguide:Refrigeration Power Consumption (arydshln-2)}
\end{table}

\IfTblCpTop{}{
\floatstyle{plain}
\restylefloat{table}
\addtolength{\abovecaptionskip}{2.5pt}
\setlength{\abovetopsep}{0pt}
}

\FloatBarrier

\subsubsection{其他候选项}
\label{sec:app:styleguide:Miscellaneous Candidates}

其他改进候选项列在本节源代码的注释中。

% Ugly line break by \co{}
%                                 __
%        atomic_store()
%
%                           seqlock_
%        t
%
%   Is there any way to prevent these breaks?
%   Maybe we need an on-the-fly script to convert such \co{}s
%   to couples of \co{}s.
%   Example:
%     \co{__atomic_store()} -> \co{__}\co{atomic_store()}
%     \co{seqlock_t} ->        \co{seqlock_}\co{t}
